\section{Markov Decision Processes}
\label{sec:mdp}

The model of \glsfirst{mdp} \cite{puterman1994markov} can be thought of as an \gls{mc} whose transition is affected by the execution of an action.
Typically, an agent is responsible for selecting this action and can therefore interact with the \gls{mdp}.
In addition, a reward function is associated with the transitions in the process and provides feedback to the agent.

Formally, the \gls{mdp} is defined as a 5-tuple $\markovdecisionprocess = (\statespace,\actionspace, \transitionfunction, \rewardfunction, \initialstatedistribution)$: $\statespace $ and $\actionspace$ are measurable sets of states and actions, respectively; $\transitionfunction(s'\vert s,a)$ is the probability of transitioning from a state $s$ to a state $s'$ by performing an action $a$; $R(s,a)$ is a random variable that counts the reward collected by the agent during the transition from state $s$ where action $a$ is applied; $\initialstatedistribution$ is the distribution of the initial state of the agent. 
We note $r(s,a)=\expectedvalue[R(s,a)]$ and make the following assumption throughout this dissertation.

\begin{ass}[Bounded reward]
    The reward is said to be bounded if, $\forall (s,a)\in\statespace\times\actionspace$,
    \begin{align*}
        0\le r(s,a)\le \maximumreward.
    \end{align*}
\end{ass}

%We will also use the notation $S_t$ or $A_t$ to refer to the random variable holding the value of the state or the action after $t$ steps in the \gls{mdp}.

\subsection{Policies}
\label{subsec:policies}
The way an agent selects its actions given its history of interaction with the environment at time $t$,  $h_t=(s_0,a_0,\dots,s_{t-1},a_{t-1},s_t)$ is called a \emph{policy}. 
If we note $\mathcal{H}_t$ the set of all histories at time $t$, then we can formally define the agent policy as a mapping from the set of histories to the set of probabilities over the action space,
\begin{align*}
    \pi:\historyspace_t\rightarrow\setofprobability\left(\actionspace\right).
\end{align*}
We follow the notation of \cite{puterman1994markov} and note the set of all policies $\Pi^{\text{HR}}$ where "$R$" stands for randomised and "$H$" for history-dependent.

Different subsets of $\Pi^{\text{HR}}$ are notable for their properties. 
The subset of Markovian policies considers policies that only depend upon the last observed state. 
This is a convenient property as it forbids the input to the policy to grow in dimension with time.
It has also been shown to contain optimal policies for the objective that we will introduce later \cite[Theorem~6.2.10 and Theorem~8.1.2]{puterman1994markov}.
%\pl{valid only for discrete state space? see https://cs.stackexchange.com/questions/93617/existence-of-optimal-policy-for-infinite-state-mdps} 
We use the superscript "$M$" to refer to Markovian policies.
When the policy is fixed--but potentially stochastic--over time, we call the policy stationary and use the superscript "$S$".
Lastly, policies can also be deterministic when the probability measure on $\actionspace$ degenerates to concentrate all the mass on a single action. 
This subset is marked with the letter "$D$".
To summarise, these sets are ordered as follows \cite[Section 2.1.5]{puterman1994markov}:
\begin{align*}
    \begin{array}{cccc}
         \Pi^{\text{SD}}&\subset\Pi^{\text{SR}}&\subset\Pi^{\text{MR}}&\subset\Pi^{\text{HR}}  \\
         \Pi^{\text{SD}}&\subset\Pi^{\text{MD}}&\subset\Pi^{\text{MR}}&\subset\Pi^{\text{HR}}\\
         \Pi^{\text{SD}}&\subset\Pi^{\text{MD}}&\subset\Pi^{\text{HD}}&\subset\Pi^{\text{HR}}.
    \end{array}
\end{align*}


\subsection{Objective}
\label{subsec:objective}
Three main objectives can be defined for an \gls{mdp}: the expected total return, the expected discounted return and the average reward \cite{puterman1994markov}.
In this thesis, we will focus on the last two types of returns.
For some \gls{mdp} $\markovdecisionprocess$, let $\probability\left(S_t=s; \pi,\initialstatedistribution\right)$ be the probability that the random variable representing the state at time $t$, $S_t$, will assume value $s$, following policy $\pi$ and given the initial state distribution $\initialstatedistribution$. 
For some policy $\policy$, we can define the expected discounted return with discount factor $\gamma$ as
\begin{align}
\label{eq:discounted_reward}
    \expectedreturn[\gamma] = \expectedvalue_{\substack{s_{t+1}\sim p(\cdot\vert s_t,a_t)\\a_t\sim\pi(\cdot\vert s_t)\\s_0\sim \initialstatedistribution}}\left[\sum_{t=1}^H \gamma^t r(s_t,a_t) \right];
\end{align}
where $H$ is the horizon that can potentially be infinite. 
Instead, the average reward performance reads, \footnote{We keep the notation $H$ instead of the more usual $T$ in order to match the notation of the discounted case and we will reserve $T$ to refer to the total number of steps, summed over episodes.}
\begin{align*}
    \averagereturnfunction = \lim_{H\rightarrow \infty}\frac{1}{H}\expectedvalue_{\substack{s_{t+1}\sim p(\cdot\vert s_t,a_t)\\a_t\sim\pi(\cdot\vert s_t)\\s_0\sim \initialstatedistribution}}\left[\sum_{t=1}^H r(s_t,a_t) \right].
\end{align*}

The expectations are taken with respect to the initial state distribution $d_0$ and the transition distribution induced by the policy.

\subsection{State-Action Occupancy Measure and Distribution}
\label{subsec:def_state_distrib}

Having some measure or distribution over the state visited by an agent is useful from a theoretical point of view. 
In fact, the performance of a policy can be directly deduced from these quantities.
We first define the state-action occupancy distribution for the discounted return criterion, 
\begin{align}
\label{def:discounted_state_distrib}
    \discountedstateoccupancydistribution[\gamma] (s,a)= (1-\gamma)\expectedvalue_{\substack{s_{t+1}\sim p(\cdot\vert s_t,a_t)\\a_t\sim\pi(\cdot\vert s_t)\\s_0\sim \initialstatedistribution}}\left[\sum_{t=1}^H \gamma^t \Ind\left(S_t=s, A_t=a\right)\right].
\end{align} 
Instead, the average reward's state-action occupancy distribution reads,
\begin{align*}
    \averagestateoccupancydistribution (s,a)= \lim_{H\rightarrow \infty}\frac{1}{H}\expectedvalue_{\substack{s_{t+1}\sim p(\cdot\vert s_t,a_t)\\a_t\sim\pi(\cdot\vert s_t)\\s_0\sim \initialstatedistribution}}\left[\sum_{t=1}^H \Ind\left(S_t=s, A_t=a\right)\right].
\end{align*}
% for the total reward criterion,
% \begin{align*}
%     \totalstateoccupancymeasure[\gamma] = \expectedvalue\left[\sum_{t=1}^\infty  \probability^\policy(s_t=s) \pi(a\vert s)\right].ò
% \end{align*}
From the state-action distributions, the state distributions can be derived as follows,
\begin{align*}
    \discountedstateoccupancydistribution[\gamma] (s)&=\int_\actionspace \discountedstateoccupancydistribution[\gamma] (s,a) \de a,
    \\
    \averagestateoccupancydistribution (s) &= \int_\actionspace\averagestateoccupancydistribution (s,a)\de a.
\end{align*}
Note that we have kept the same notation, but the input to the function avoids confusion. 
To further simplify the notation, we will sometimes drop the subscript ``$\gamma$'' or ``AVG'' when no confusion about the objective can be made.

Finally, we define the state-action occupancy measure \cite{laroche2022non} in the discounted case,
\begin{align*}
    \discountedstateoccupancymeasure[\gamma] (s,a) &= \expectedvalue_{\substack{s_{t+1}\sim p(\cdot\vert s_t,a_t)\\a_t\sim\pi(\cdot\vert s_t)\\s_0\sim \initialstatedistribution}}\left[\sum_{t=1}^H \gamma^t \Ind\left(S_t=s, A_t=a\right)\right]
    \\
    &= \frac{\discountedstateoccupancydistribution[\gamma] (s,a)}{1-\gamma},
\end{align*}
and for the average reward,
\begin{align*}
    \averagestateoccupancymeasure[\pi] (s,a) = \lim_{H\rightarrow \infty}\expectedvalue_{\substack{s_{t+1}\sim p(\cdot\vert s_t,a_t)\\a_t\sim\pi(\cdot\vert s_t)\\s_0\sim \initialstatedistribution}}\left[\sum_{t=1}^H \Ind\left(S_t=s, A_t=a\right)\right].
\end{align*}
The discounted return of a policy can be deduced from the state-action occupancy measure as follows \cite[Lemma~1]{laroche2022non},
\begin{align*}
    \expectedreturn[\gamma] &= \int_{\statespace,\actionspace} r(s,a)d\discountedstateoccupancymeasure[\gamma]  (\de s,\de a).
\end{align*}

\subsection{Value Functions}
An important concept for quantifying the performance of a policy in \gls{rl} is that of \emph{value functions}. 
In the discounted case, the state-action value function is the expected discounted sum of rewards that the agent will collect starting in some state $s$, applying the action $a$ first and following the policy $\pi$ thereafter. 
It reads,
\begin{align}
    Q^\pi(s,a) = \expectedvalue_{\substack{s_{t+1}\sim p(\cdot\vert s_t,a_t)\\a_t\sim\pi(\cdot\vert s_t)}}
    \left[\left.\sum_{t=0}^H \gamma^t r(s_t,a_t)\right\vert s_0=s, a_0=a\right].\label{eq:state_action_value_function}
\end{align}
This function is called \emph{Q-function}.
Naturally, it is also possible to define a state value function or \emph{value function} for short. 
It is the expected discounted sum of reward that the agent will collect from starting in some state $s$ and immediately following the policy $\pi$.
It is defined from the Q-function as
\begin{align*}
    V^\pi(s)=\expectedvalue_{a\sim\pi(\cdot\vert s)}[Q^\pi(s,a)].
\end{align*}

A useful concept to compare the effect of an action is the advantage function $\advantagefunction$.
It quantifies the advantage obtained by deviating from a given policy only for the next step. 
Formally,
\begin{align}
\label{eq:advantage_function}
    \advantagefunction(s,a) = Q^\pi(s,a) - V^\pi(s)
\end{align}

In \gls{rl}, we are interested in the maximum value that the state-action and state-value functions can reach. 
The optimal state and state-action value functions are 
\begin{align*}
    \optimalstatevaluefunction (s)&= \sup_{\pi\in\Pi^{\text{HR}}}V^\pi(s),
    \\
    \optimalstateactionvaluefunction (s,a)&= \sup_{\pi\in\Pi^{\text{HR}}}Q^\pi(s,a).
\end{align*}
The policy $\pi\in\Pi^{\text{HR}}$ that achieves these values is noted $\optimalpolicy$. 
Note that there exists a policy in $\Pi^{\text{SR}}$ which is optimal  \cite[Theorem~6.2.10]{puterman1994markov}. 
Therefore, we can restrict the search for an optimal policy accordingly. 

\subsection{Bellman Equations}
A core concept for finding an optimal value function in \glspl{mdp} is \emph{dynamic programming} \cite{bellman1954theory}. 
Its idea is to break down a problem into smaller problems.
Applied to \gls{rl}, the problem of computing the value function at any state is decomposed into recursively finding the value function at each step, as a function of future value functions. 
This is possible thanks to the \emph{Bellman expectation operators} \cite{bellman1966dynamic}.

\begin{defi}[Bellman Expectation Operators]
    Consider an \gls{mdp} and a policy $\pi\in\Pi^{\text{SR}}$ and let $\mathcal{B}(\statespace)$ be the set of bounded measurable on $\statespace$. 
    The Bellman expectation operator for the state value function $\bellmanoperator_V:\mathcal{B}(\statespace)\mapsto\mathcal{B}(\statespace)$ is defined for $f\in\mathcal{B}$ and $s\in\statespace$ as follows,
    \begin{align}
    \label{eq:bellman_operator_V}
        (\bellmanoperator_V f)(s) = \int_{\actionspace}\pi(\de a\vert s)\left(r(s,a) + \gamma  \int_\statespace  \transitionfunction(\de s'\vert s,a) f(s')\right).
    \end{align}
    Similarly, the Bellman expectation operator for the state-action value function $\bellmanoperator_Q:\mathcal{B}(\statespace\times\actionspace)\mapsto\mathcal{B}(\statespace\times\actionspace)$ is defined for $f\in\mathcal{B}$ and $s\in\statespace$, $a\in\actionspace$ as follows,
    \begin{align}
    \label{eq:bellman_operator_Q}
        (\bellmanoperator_Q f)(s,a) = r(s,a) + \gamma  \int_\statespace  \transitionfunction(\de s'\vert s,a)\int_\actionspace \pi(\de a'\vert s')f(s',a').
    \end{align}
\end{defi} 

These operators are $\gamma$-contractions for the $L_\infty$-norm \cite[Proposition~6.2.5]{puterman1994markov} if $\gamma<1$.
Recall that a $\gamma$-contraction in the $L_\infty$-norm is an operator $T:X\mapsto X$ and $f,g\in X$, $T$ such that,
\begin{align*}
    \lVert Tf - Tg\rVert_\infty\le\gamma\lVert f-g\rVert_\infty.
\end{align*}
These properties of the Bellman expectation operators place them under the conditions of the Banach fixed-point theorem. 
It implies that their recursive application converges to a unique fixed point.
The Q-function is the fixed-point for \Cref{eq:bellman_operator_Q} and the state value function for \Cref{eq:bellman_operator_V}.
From there, the \emph{Bellman expectation equations} are defined:
\begin{align}
    Q^\pi &= r(s,a) + \gamma  \int_\statespace  \transitionfunction(\de s'\vert s,a)V^\pi(s'),
    \label{eq:bellman_equation_Q}\\
    V^\pi &= \int_\actionspace \pi(\de a\vert s) Q^\pi(s,a).
    \label{eq:bellman_equation_V}
\end{align}

Interestingly, the optimal state-action and state value function are also fixed points of similar operators called \emph{Bellman optimality operators} and defined hereafter. 

\begin{defi}[Bellman Optimality Operators]
    Consider an \gls{mdp} and let $\mathcal{B}(\statespace)$ be the set of bounded measurable functions on $\statespace$. 
    The Bellman optimality operator for the state value function $\optimalbellmanoperator_V:\mathcal{B}(\statespace)\mapsto\mathcal{B}(\statespace)$ is defined for $f\in\mathcal{B}$ and $s\in\statespace$ as follows,
    \begin{align}
    \label{eq:bellman_optimal_operator_V}
        (\optimalbellmanoperator_V f)(s) = \sup_{a\in\actionspace}\left\{r(s,a) + \gamma  \int_\statespace  \transitionfunction(\de s'\vert s,a) f(s')\right\}.
    \end{align}
    Similarly, the Bellman expectation operator for the state-action value function $\optimalbellmanoperator_Q:\mathcal{B}(\statespace\times\actionspace)\mapsto\mathcal{B}(\statespace\times\actionspace)$ is defined for $f\in\mathcal{B}$ and $s\in\statespace$, $a\in\actionspace$ as follows,
    \begin{align}
    \label{eq:bellman_optimal_operator_Q}
        (\optimalbellmanoperator_Q f)(s,a) = r(s,a) + \gamma  \int_\statespace  \transitionfunction(\de s'\vert s,a)\sup_{a'\in\actionspace} \left\{f(s',a')\right\}.
    \end{align}
\end{defi} 

These operators are also $\gamma$-contractions \cite[Proposition~6.2.4]{puterman1994markov} and respectively admit the optimal state value function and the optimal state-action value function as fixed-points from the Banach fixed-point theorem \cite[Proposition~6.2.5]{puterman1994markov}. 

Similarly, the optimal value functions satisfy the following \emph{Bellman optimality equations}:
\begin{align}
    Q^\star(s,a) &= r(s,a) + \gamma  \int_\statespace  \transitionfunction(\de s'\vert s,a)V^\star(s'),
    \label{eq:bellman_optimal_equation_Q}\\
    V^\star(s) &= \sup_{a\in\actionspace} Q^\star(s,a).
    \label{eq:bellman_optimal_equation_V}
\end{align}


\subsection{Smooth Markov Decision Process}
\label{subsec:smooth_mdp}

Considering the whole class of \glspl{mdp} implies considering some very intricate problems that can be difficult to solve efficiently.
Making assumptions to obtain a restricted set of \glspl{mdp} can drastically reduce the complexity while maintaining the realism of the model. 
One such assumption that is ubiquitous in \gls{rl} is about the \emph{smoothness} of the \gls{mdp}.
To define what is intended by smoothness, we first recall the concept of Lipschitzness.  

\begin{defi}[Lipschitz continuity]
    Consider $(X,\distance_X)$ and $(Y,\distance_Y)$ two metric spaces and $f$ a function $f:X \rightarrow Y$. Then, $f$ is $L$-Lipschitz continuous ($L$-LC) for $L>0$ if 
    \begin{align*}
        \forall x,x'\in X,\; \distance_Y(f(x),f(x'))\leq L \distance_X(x,x').
    \end{align*}
\end{defi}

Throughout this dissertation, for sets $X$ such that $X\subset \mathbb{R}^n$, we consider the Euclidean distance defined for $x,x'\in X$ as $\distance_{X}(x,x')=\lVert x-x'\rVert_2$.
For probability distributions, we will consider the $L_1$-Wasserstein distance\footnote{By abuse of language and for short, we will refer to the $L_1$-Wasserstein distance as the Wasserstein distance in this thesis.} which we define hereafter.

\begin{defi}[$L_1$-Wasserstein distance, \cite{villani2009optimal}]
    Let $\mu,\nu$ be two probabilities with sample space $\Omega$ is:
    \begin{align*}
        \wassersteindistance[1](\mu\Vert \nu)=\sup_{\left\Vert f\right\Vert_L\leq 1} \left\vert \int_{\Omega} f(\omega)(\mu-\nu)\; (\de \omega) \right\vert,
    \end{align*}
    where $\left\Vert f\right\Vert_L$ is the Lipschitz semi-norm of $f$:
    \begin{align*}
        \left\Vert f\right\Vert_L=\sup_{x,x'\in X, x\neq x'} \frac{\distance_Y(f(x),f(x'))}{\distance_X(x,x')}.
    \end{align*}
\end{defi}

The Lipschitzness can be understood as some notion of smoothness as it limits the speed of growth of a function.
We are now ready to describe the smoothness of an \gls{mdp}.

\begin{defi}[Lipschitz \gls{mdp}, \cite{rachelson2010locality} ]
\label{def:lip_mdp}
An \gls{mdp} is $(L_P,L_r)$-LC if, $\forall(s,a),(s',a')\in\statespace\times\actionspace$
\begin{align*}
    & \wassersteindistance[1](\transitionfunction(\cdot\vert s,a)\Vert \transitionfunction(\cdot\vert s',a'))\leq L_P \left( \distance_{\statespace}(s,s')  + \distance_{\actionspace}(a,a')\right),
    \\
    & \left\vert r(s,a)-r(s',a') \right\vert \leq L_r \left( \distance_{\statespace}(s,s') + \distance_{\actionspace}(a,a')\right).
\end{align*}
\end{defi}

A similar assumption of smoothness can be made for the policy, which is given in the definition.

\begin{defi}[Lipschitz policy]
\label{def:lip_policy}
A stationary Markovian policy $\pi$ is $L_\pi$-LC if, $\forall s,s'\in\statespace$ 
\begin{align*}
    \wassersteindistance[1](\pi(\cdot\vert s)\Vert \pi(\cdot\vert s'))\leq L_\pi \distance_{\statespace}(s,s').
\end{align*}
\end{defi}


These assumptions have great implications. 
For example, the identification of dominating actions is made simpler by guaranteeing that this action dominates over a ball in the state space \cite{rachelson2010locality}. 
The same reference also provides a result on the Lipschitzness of the Q-function, which is recalled below.

\begin{thm}[Lipschitzness of the Q-function, Theorem~1 of \cite{rachelson2010locality}]
\label{th:lip_Q_function}
    Consider an $(L_P,L_r)$-LC \gls{mdp} and an $L_\pi$-LC policy $\policy$. If $\gamma L_P(1+L_\policy)\le 1$, then $Q^\pi$ is $L_Q$-LC with
    \begin{align*}
        L_Q=\frac{L_r}{1-\gamma L_P(1+L_\policy)}.
    \end{align*}
\end{thm}

Finally, a more recent notion of Lipschitzness with respect to time has been introduced in \cite[Assumption 4.1]{metelli2020control}. 
The assumption will be particularly useful in the case of delays, as the smoothness of the trajectories is a key factor for the agent to predict the effect of its actions. 

\begin{defi}[Time-Lipschitz MDP, Assumption~4.1 of \cite{metelli2020control}]
\label{def:time_lip}
An \gls{mdp} is $L_T$-Time Lipschitz Continuous ($L_T$-TLC) if, $\forall (s,a)\in \statespace\times\actionspace$
\begin{align*}
    \wassersteindistance[1](\transitionfunction(\cdot\vert s,a)\Vert \delta_s)\le L_T,
\end{align*}
where $\delta_s$ is defined from the Dirac measure $\delta$ as $\delta_s(x)=\delta(x-s)$.
\end{defi}

It can seem that this definition is not consistent with the traditional definition of Lipschitzness; however, as we will demonstrate, it can be seen as Lipschitzness w.r.t. the number of steps elapsed.
Indeed, the r.h.s. can be read as $L_T\cdot 1$ where $1$ corresponds to the unit of time, the step, of the \gls{mdp}. 
Second, for the l.h.s., we could rewrite $p(\cdot\vert s,a)$ as $p^1(\cdot\vert s,a)$ to emphasise that the transition happens in a step.
Considering the application of more than one action, for instance, the sequence $\boldsymbol a=(a_1,\dots,a_n)$ of $n$ actions, we could introduce the notation $p^n(\cdot\vert s,\boldsymbol a)$.
It would represent the sequential application of the actions contained in $\boldsymbol a$ starting from state $s$.
Finally, we could extend the notation to include $p^0(\cdot\vert s)=\delta_s$ the effect of applying no action at state $s$.
With this new notation, the equation inside the TLC definition becomes:
\begin{align*}
    \wassersteindistance[1](p^1(\cdot\vert s,a)\Vert p^0(\cdot\vert s))\le L_T\cdot (1-0).
\end{align*}
It becomes clearer that the Lipschitzness is with respect to the step of the \gls{mdp}.
We will prove an extension of this result for the distance between $p_{\boldsymbol a}^n(\cdot\vert s)$ and $p^0(\cdot\vert s)$ in \Cref{pp:tlc_delay}.


\subsection{POMDP}
\label{subsec:pomdp}

In this subsection, we briefly introduce the concept of \gls{pomdp} \cite{kaelbling1998planning} which has similarities with the delayed processes, as we will see later.
\gls{pomdp} extend the framework of \gls{mdp} to consider situations where the agent can no longer observe the current state but only has access to partial information on it. 
More formally, a \gls{pomdp} is a tuple $(\statespace,\actionspace, \transitionfunction, \rewardfunction, \initialstatedistribution, \Omega, O)$ where, compared to an \gls{mdp}, two elements are added. 
First, $\Omega$ is the set of observations that the agent may receive instead of the states themselves. 
Second, $O:\statespace\times\actionspace\mapsto \setofprobability(\Omega)$ is called the observation function and defines a probability distribution over $\Omega$. 
Specifically, if an agent arrives at state $s'$ by taking action $a$, then $O(o; s',a)$ is the probability that the agent will receive $o$ as observation.
The reward and transition functions still depend on the current state of the environment, but this state is unknown to the agent. 
The objective remains the one defined in \Cref{subsec:objective}.

For solving \gls{pomdp}, a traditional approach is to keep track of the probability distribution of the current state, called belief, by updating this probability with each new observation collected. 
Doing so, one can cast the problem back to an \gls{mdp} where the state would be the aforementioned belief.
A similar concept is also found in the control theory literature \cite{kumar2015stochastic} under the name of information state, and, as we will see in \Cref{sec:related_works}, it has also been used for the delayed problem. 
For more information on \gls{pomdp}, we refer the reader to \cite{hauskrecht2000value, spaan2012partially}.
